% =====================================================================
% HMM: вывод Forward-формулы
% =====================================================================
\begin{frame}{HMM: вывод Forward-формулы (Rabiner 1989)}

\textbf{Определим} $\alpha_t(i) = P(O_{1:t},\, S_t = i \mid \lambda)$
--- вероятность, что мы наблюдали $o_1, \ldots, o_t$ и сейчас находимся
в состоянии $i$.

\bigskip

\textbf{Инициализация} (по определению совместного):
\[
\alpha_1(i) \;=\; \pi_i\, b_i(o_1),
\quad i \in \{1, \ldots, N\},
\]
где $\pi_i = P(S_1 = i)$ --- начальное распределение,
$b_i(o) = P(o \mid S=i)$ --- функция эмиссии.

\bigskip

\textbf{Индукция $t \to t+1$:} маргинализуем по предыдущему состоянию
$S_t = j$:
\[
\alpha_{t+1}(i)
\;=\; P(O_{1:t+1},\, S_{t+1}=i)
\;=\; \sum_{j=1}^{N} P(O_{1:t+1},\, S_t = j,\, S_{t+1} = i).
\]

\end{frame}

\begin{frame}{HMM: вывод Forward-формулы (продолжение)}

\textbf{Цепной разбор} по правилу произведения и марковскому свойству:
\[
P(O_{1:t+1}, S_t=j, S_{t+1}=i)
= \underbrace{P(O_{1:t}, S_t=j)}_{=\alpha_t(j)}\;
  \underbrace{P(S_{t+1}=i \mid S_t=j)}_{=a_{j,i}}\;
  \underbrace{P(o_{t+1} \mid S_{t+1}=i)}_{=b_i(o_{t+1})}.
\]

\bigskip

Подставляя в маргинализацию, получаем рекурсивную формулу:

\bigskip

\[
\boxed{\;
\alpha_{t+1}(i) \;=\; b_i(o_{t+1}) \sum_{j=1}^{N} \alpha_t(j)\, a_{j,i}\;}
\]

\bigskip

\textbf{Сложность шага:} $O(N^2)$ (для каждого из $N$ новых состояний
суммируем по $N$ предыдущим). Для нашей задачи переходы
\emph{разреженные} (только stay/advance/skip), поэтому шаг $O(N)$.

\end{frame}

\begin{frame}{HMM: предсказание и confidence}

После каждого события считаем $\alpha_t(i)$ для всех $i$, затем:

\bigskip

\textbf{Предсказание позиции:}
\[
\hat\imath_t \;=\; \argmax_{i} \alpha_t(i).
\]

\bigskip

\textbf{Уверенность} (confidence): мера «остроты» распределения
$\alpha_t$:
\[
\mathrm{conf}_t \;=\; \frac{\alpha_t(\hat\imath_t)}{\sum_j \alpha_t(j)}.
\]

\bigskip

\textbf{Зачем это нужно для нас:}
\begin{itemize}
\item HMM \emph{выдаёт распределение}, а не одно число (в отличие от OLTW).
\item Это распределение можно \rlhi{подать на вход RL-агенту} как
      признак неопределённости.
\item Confidence используется как переключатель в гибридной модели:
      $\mathrm{conf}_t > \theta \Rightarrow$ HSMM,
      иначе $\Rightarrow$ OLTW.
\end{itemize}

\bigskip

\textbf{Численная стабильность:} в realtime мы нормируем $\alpha_t$ на
сумму после каждого шага --- иначе они уходят в \texttt{underflow}.

\end{frame}
